\chapter{Theoretical Background}
\label{ch:theory}

\section{Introduction}
\label{sec:theory-intro}

The design and implementation of ScholarPath AdDU is grounded in nine established theoretical frameworks drawn from computer science, information science, and software engineering. Together, these theories provide formal justification for every major architectural decision, from database structure and document management to security, user experience, and development methodology. This chapter situates each theory within the overarching Information Systems Success Model \cite{ref20} adopted in Chapter~\ref{ch:rrl}. Relational Database Theory and Metadata-Driven Document Management underpin System Quality by ensuring structural integrity and storage efficiency. Information Foraging Theory and Event-Driven Architecture support Information and Service Quality through improved discoverability of applicant data and proactive communication. Rule-Based Screening Theory, bounded by a human-in-the-loop decision rule, explains how the platform turns the SOP and the five school rule sets into flags that inform, but never make, committee decisions. Finally, Cloud Infrastructure Theory, RBAC, Empirical Process Control, and the HCI/ISO 25010 framework collectively ensure scalability, legal compliance, iterative development, and empirically validated usability.

\section{Relational Database Theory and Normalization}
\label{sec:relational}

Consolidating applicant academic records, household financial documents, interview scores, five school rule sets, and an audit log of human decisions into a consistent repository requires a rigorous structural foundation. Relational Database Theory, specifically the normalization principles formalized by Codd \cite{ref17}, provides this foundation by organizing data into interconnected two-dimensional tables governed by first-order predicate logic. Within the DeLone and McLean \cite{ref20} framework, a normalized 3NF schema is the foundational enabler of System Quality, specifically the sub-dimension of data integrity, as it eliminates the inconsistencies that would otherwise corrupt the data presented to committees.

\subsection{The Relational Model and Anomaly Prevention}
\label{sec:anomalies}

Codd \cite{ref17} defined normalization as the systematic process of organizing relations to reduce redundancy and prevent modification anomalies. In the AdDU context, three anomaly types directly threaten the integrity of an unnormalized scholarship admission database. Table~\ref{tab:anomalies} outlines each anomaly and its resolution within ScholarPath AdDU.

\begin{table}[htbp]
  \centering
  \small
  \caption{Database Anomalies and Their Resolution in ScholarPath AdDU}
  \label{tab:anomalies}
  \begin{tabular}{@{}p{2.4cm}p{5.2cm}p{6.8cm}@{}}
    \toprule
    Anomaly Type & Definition & Resolution in ScholarPath AdDU \\
    \midrule
    Insertion Anomaly & Facts cannot be recorded because they lack a required primary key from another entity. & School rule sets and track definitions are stored independently of applications, so a subcommittee's criteria for a new cycle can be configured before any applicant registers. \\
    Update Anomaly & A single logical fact recorded across multiple rows creates contradictions when updates are partially executed. & Each baseline threshold, phase window, and school criterion exists in one canonical row; an update propagates to every partition and concurrent committee query. \\
    Deletion Anomaly & Deleting one record unintentionally destroys data representing an entirely different, independent fact. & Withdrawing an application does not remove its audit log entries or the school's rule set, which are stored in independent relations. \\
    \bottomrule
  \end{tabular}
\end{table}

\subsection{Progression Through Normal Forms}
\label{sec:normal-forms}

The ScholarPath AdDU schema adheres to the Third Normal Form (3NF), progressing through Codd's sequential requirements. First Normal Form (1NF) mandates atomic, scalar values with no repeating groups; prerequisite documents such as the application form, report card, and BIR Form 2316 are therefore stored as discrete, queryable records rather than comma-separated strings. Second Normal Form (2NF) eliminates partial dependencies, ensuring school and track details are stored independently of application records. Third Normal Form (3NF) eliminates transitive dependencies, ensuring, as Kent \cite{ref32} summarizes, that every non-key attribute provides ``a fact about the key, the whole key, and nothing but the key.'' Adhering to 3NF guarantees that any update to a configured parameter, such as a revised baseline average or an extended Phase~2 window, propagates from a single authoritative record to all related queries.

\section{Centralized Information Systems and Cloud Infrastructure Theory}
\label{sec:cloud-theory}

The decision to deploy ScholarPath AdDU as a single unified web platform, rather than distributing applicant data across the files of five school subcommittees, is grounded in centralized IT infrastructure theory \cite{ref34}. Centralized infrastructure consolidates computing resources and data into one managed environment, enabling uniform policy enforcement and the elimination of data silos that cause policy desynchronization across departments \cite{ref41}. By hosting the entire SOP-aligned Application Lifecycle Path on a single Supabase backend, changes to submission windows or school criteria are immediately reflected for all users. Modern cloud platforms mitigate the traditional risk of centralization, a single point of failure, through hardware virtualization and physical redundancy. Mell and Grance \cite{ref34} define cloud computing as on-demand access to a shared pool of configurable computing resources characterized by rapid elasticity and measured service. For ScholarPath AdDU, Supabase's managed PostgreSQL infrastructure provides high availability, and the platform's audit log of every status change supports both compliance with the Philippine Data Privacy Act of 2012 (RA 10173) and the traceability expected in institutional quality audits \cite{ref52}. % TODO(OPEN-1): name the ISO standard once confirmed.
In the IS Success Model, centralized cloud infrastructure directly supports System Quality through 24/7 availability and Service Quality through uniform, institution-wide policy enforcement.

\section{Metadata-Driven Document Management}
\label{sec:metadata}

Manual, paper-based document workflows create severe operational bottlenecks for the University Scholarship Office, which must receive and check a complete physical file from every applicant. ScholarPath AdDU resolves this through metadata-driven document management, which decouples physical file storage from logical document associations \cite{ref25}. Table~\ref{tab:storage} contrasts traditional hierarchical storage with the metadata-driven approach.

\begin{table}[htbp]
  \centering
  \small
  \caption{Traditional Hierarchical Storage versus Metadata-Driven Architecture}
  \label{tab:storage}
  \begin{tabular}{@{}p{2.8cm}p{5.4cm}p{6.2cm}@{}}
    \toprule
    Feature & Traditional Hierarchical Storage & Metadata-Driven Architecture \\
    \midrule
    Location Strategy & Files are locked into a rigid folder path. & Files are stored in a flat virtual pool; organization is determined dynamically by relational metadata tags. \\
    Submission Phases & A single complete file must be assembled and submitted at once. & Each document is tagged with its submission phase, so Phase~1 and Phase~2 documents accumulate against one application. \\
    Corrections & A replaced page overwrites or is filed separately from the original. & A corrected upload during a Conditional Revert is stored as a new version of the same requirement, preserving the submission history. \\
    Retrieval Mechanism & Staff navigate a folder hierarchy to locate files. & Staff query by document type, phase, upload date, or verification status. \\
    \bottomrule
  \end{tabular}
\end{table}

ScholarPath AdDU operationalizes this theory as the Document Vault. When an applicant uploads a prerequisite file, such as a report card or BIR Form 2316, it is stored exactly once in a Supabase Storage bucket and assigned a unique URI. The Applications table references that file through foreign keys, and its metadata records the submission phase and verification status. University Scholarship Office staff can therefore query and verify documents through metadata tags rather than folder hierarchies.

\section{Rule-Based Screening and Human-in-the-Loop Decision Support}
\label{sec:rule-screening}

Rule-Based Expert Systems (RBS) theory represents domain expertise as conditional IF-THEN logic \cite{ref15}. ScholarPath AdDU applies this theory in a deliberately bounded way. The rules encode the SOP baseline and the five School Scholarship Subcommittee rule sets, but their conclusions are flags, sort orders, and alerts, never verdicts on an applicant. The decision that changes an applicant's status is always made and recorded by a person. This boundary answers the interpretability concern that Kanwal et al. \cite{ref50} raise about automated scholarship eligibility decisions, and it is the architectural expression of the institutional directive that no algorithm selects scholars \cite{ref52}. % TODO(L-6): add team-supplied decision-support / human-oversight sources as ref53+.
In the DeLone and McLean \cite{ref20} model, rule-based screening supports Information Quality, specifically the completeness and relevance of the data a committee sees, while the human decision rule supports the accountability of the Net Benefits the system produces.

\subsection{Knowledge Representation and the Inference Engine}
\label{sec:knowledge}

Buchanan and Duda \cite{ref15} describe expert systems as consisting of three interdependent components: a Knowledge Base of facts and IF-THEN rules, a Working Memory storing the current input state, and an Inference Engine that applies rules to derive conclusions via the logical principle of modus ponens. In ScholarPath AdDU, the Knowledge Base is the SOP baseline (an entrance exam ``ACCEPTED'' remark, a high school general average of at least 85\%, and the required documents) together with the five school rule sets; the Working Memory is an application's verified data at a given lifecycle stage; and the Inference Engine applies a forward-chaining strategy from those data to a set of screening flags. The conclusions of the engine are routed to a staff review queue or displayed in the Committee Portal, where a person acts on them.

\subsection{Rule Precedence and Committee Resolution}
\label{sec:precedence}

A core design challenge in rule-based systems is conflict resolution: determining behavior when multiple rules apply to the same data state. In ScholarPath AdDU, conflicts arise when a track rule and a school rule point in different directions. For example, a certified Rank~1 graduate qualifies for the fixed Jubilee grant rate, while the SEA rule set applies a hard entrance exam cutoff to the same applicant's program choice. The platform does not resolve such conflicts automatically. It raises both flags, shows them together to the School Scholarship Subcommittee, and records the subcommittee's decision. % TODO(OPEN-3): Jubilee guarantee vs. school hard cutoff.
Table~\ref{tab:conflict} maps the conflict resolution strategies to their system context.

\begin{table}[htbp]
  \centering
  \small
  \caption{Rule Conflict Handling in ScholarPath AdDU Screening}
  \label{tab:conflict}
  \begin{tabular}{@{}p{3.0cm}p{5.0cm}p{6.4cm}@{}}
    \toprule
    Strategy & Theoretical Definition & ScholarPath AdDU Implementation \\
    \midrule
    Specificity Priority & Rules with more specific conditions are evaluated before general rules. & School-specific rules, such as the SON zero-conditional policy, are evaluated and displayed alongside the general SOP baseline flags. \\
    Flag, Do Not Fire & Conflicting conclusions are surfaced rather than acted upon. & When a track rule and a school rule conflict, both flags are shown to the subcommittee; no status change occurs until a member records a decision. \\
    Human Resolution & A designated authority resolves conflicts that rules cannot. & The School Scholarship Subcommittee's recorded decision stands and is written to the audit log with the acting member and timestamp. \\
    \bottomrule
  \end{tabular}
\end{table}

\section{Information Foraging Theory and Committee Filtering}
\label{sec:ift}

The cognitive processes through which users navigate complex digital environments are formalized by Information Foraging Theory (IFT), developed by Pirolli and Card \cite{ref37} at PARC. IFT applies optimal foraging theory from evolutionary ecology to model how users maximize the rate of information gain while minimizing cognitive cost. IFT rests on three core constructs: Information Patches (localized clusters of related content within an interface), Information Diet (the specific subset of sources a user pursues to satisfy a goal), and Information Scent (the user's cognitive estimate of how likely a given path is to yield relevant results). When perceived scent drops below a threshold, users abandon the current patch and navigate elsewhere \cite{ref37}. In the AdDU context, a subcommittee member facing an undifferentiated list of approximately 1,326 applications experiences weak information scent and high foraging cost. ScholarPath AdDU addresses this by partitioning applications by school, so each subcommittee begins in its own patch, and by exposing the underlying metadata as filters, such as track, program choice, verification status, and interview score, that the member intersects to refine the patch \cite{ref4,ref43}. Backend compound indexing and a 300-millisecond client-side debounce ensure fast query response, preserving the low temporal cost critical to effective information foraging.

\section{Event-Driven Architecture for Asynchronous Notification}
\label{sec:eda}

Missed submission windows and uncorrected files are attributable in part to passive, manual communication channels. ScholarPath AdDU addresses this through Event-Driven Architecture (EDA), which shifts the notification system from a synchronous request-response model to an asynchronous, event-triggered one \cite{ref28}. EDA is built around the production, detection, and consumption of events, defined as significant changes in system state, and operates on the principle of loose coupling, where components interact independently with no awareness of each other's internal logic. By adopting EDA, the notification subsystem operates without blocking the main application server. Table~\ref{tab:eda} maps the event triggers to their corresponding EDA components.

\begin{table}[htbp]
  \centering
  \small
  \caption{Event-Driven Architecture Trigger Mapping for the Automated Notification Subsystem}
  \label{tab:eda}
  \begin{tabular}{@{}p{2.2cm}p{4.0cm}p{4.0cm}p{4.2cm}@{}}
    \toprule
    EDA Component & Temporal Trigger (Deadlines and Revert Clock) & Transactional Trigger (Human Status Change) & Transactional Trigger (Release and Vacancy) \\
    \midrule
    Event Producer & Scheduled Cron job checking phase windows and five-day Revert clocks daily & Staff or subcommittee member committing a status change & Office of Admission recording the release of results; an awardee declining an offer \\
    Event Router & Database logic identifying applications near a deadline or past the Revert window & Database webhook detecting the change upon commit & Database webhook detecting the release or the freed slot \\
    Event Consumer & Edge Function sending an SMS reminder via Twilio, or flagging the file Lapsed and alerting staff & Edge Function sending a Conditional Revert notice via SendGrid and Twilio & Edge Function sending outcome notices to applicants, or alerting the subcommittee with a proposed waitlisted candidate \\
    \bottomrule
  \end{tabular}
\end{table}

This decoupled architecture ensures that third-party API latency, inherent in calls to Twilio and SendGrid, is absorbed entirely by the Edge Function layer, allowing the core web server to scale efficiently under concurrent load. It also enforces the ordering required by the SOP: an applicant's outcome notice is produced only by the release event of the Office of Admission, never by the subcommittee's decision itself.

\section{Role-Based Access Control and Data Privacy}
\label{sec:rbac-theory}

Scholarship admission systems process highly sensitive personal and socioeconomic data, necessitating rigorous access control. ScholarPath AdDU implements the NIST Role-Based Access Control (RBAC) model, formalized by Ferraiolo and Kuhn \cite{ref24} and standardized by NIST in 2004. RBAC replaced error-prone, per-user Access Control Lists (ACLs) with a tripartite structure: Users are assigned Roles representing their organizational function, and Permissions are attached exclusively to those Roles \cite{ref38}. This structure inherently enforces two foundational security principles: the Principle of Least Privilege, which limits access to the minimum required for each user's duties, and Separation of Duties, which distributes critical workflows across roles to prevent unilateral fraud or error. Separation of Duties is especially visible in ScholarPath AdDU: University Scholarship Office staff verify documents, interview panels score, School Scholarship Subcommittees select, and the Office of Admission releases results, so no single role carries an application from submission to award. RBAC is also the direct mechanism by which the platform complies with RA 10173. Applicants hold self-access only; staff hold access to all applications for verification; interview panels see only their assigned applications; each subcommittee sees only its own school's partition; and the System Administrator configures the platform without authority over any applicant's status (Table~\ref{tab:rbac}). These boundaries are enforced at the database engine level through PostgreSQL Row Level Security (RLS) policies. Within the IS Success Model, RBAC is the mechanism by which the platform satisfies the Service Quality dimension, specifically the assurance sub-dimension.

\section{Empirical Process Control Theory in Agile Development}
\label{sec:agile-theory}

The development of ScholarPath AdDU is governed by Empirical Process Control theory, the foundational pillar of Agile methodology \cite{ref39}. Unlike the Waterfall model, which assumes complete upfront specification of all requirements, empirical process control acknowledges that complex software systems operate in volatile, ambiguous environments where requirements emerge through iterative stakeholder interaction. The theory rests on three interdependent pillars \cite{ref39}: Transparency, requiring that all development progress and system artifacts be visible and understood by both developers and stakeholders; Inspection, requiring frequent examination of the system against standard metrics to detect variances; and Adaptation, requiring immediate process adjustment when inspection reveals unacceptable deviation. The five school rule sets and the SOP contain procedural details that cannot be fully specified before development begins, several of which remain under confirmation with the University Scholarship Office. The development of ScholarPath AdDU has already exercised all three pillars. Key Informant Interviews with the admissions office served as Inspection, revealing that an earlier automated matching design conflicted with institutional governance, and the redesign into a data-feeding architecture was the resulting Adaptation \cite{ref52}. The Agile SDLC described in Chapter~\ref{ch:methodology} continues this cycle through alpha testing and User Acceptance Testing (UAT). In the context of the IS Success Model, empirical process control supports System Quality by enabling continuous verification that the system's functional logic aligns with stakeholder requirements. Hoda et al. \cite{ref27} conducted a systematic mapping of agile practices across 1,540 studies, confirming that transparency, iterative inspection, and stakeholder-driven adaptation remain the dominant predictors of project success in environments characterized by evolving or partially specified requirements, conditions that describe AdDU's scholarship admission procedures.

\section{Human-Computer Interaction and the ISO/IEC 25010 Model}
\label{sec:hci}

While robust backend architecture ensures computational reliability, the practical success of an information system is ultimately determined by user adoption. The interface design and formal evaluation of ScholarPath AdDU are therefore guided by Human-Computer Interaction (HCI) theory and the ISO/IEC 25010:2023 Software Quality Evaluation standard \cite{ref23,ref31}. HCI focuses on designing interactive computing systems that align with users' cognitive capacities and actual workflows \cite{ref23}. For applicants, the legacy process of assembling a complete physical file, submitting it in person, and waiting without status updates imposes unnecessary cognitive and logistical overhead. ScholarPath AdDU alleviates this through an applicant dashboard showing the current lifecycle stage, outstanding documents, and correction requests, and through a mobile-responsive architecture enabling uploads from any device. For staff and subcommittee members, the Committee Portal replaces manual collation with partitioned, filterable views.

The ISO/IEC 25010:2023 SQuaRE standard provides the formal evaluation framework, defining product quality across eight measurable characteristics. The ScholarPath AdDU prototype assessment operationalizes two dimensions, Functional Suitability and Performance Efficiency, as detailed in Table~\ref{tab:iso25010}.

\begin{table}[htbp]
  \centering
  \small
  \caption{ISO/IEC 25010:2023 Evaluation Dimensions for ScholarPath AdDU}
  \label{tab:iso25010}
  \begin{tabular}{@{}p{2.8cm}p{3.2cm}p{8.4cm}@{}}
    \toprule
    Characteristic & Sub-Characteristic & Evaluation Metric for ScholarPath AdDU \\
    \midrule
    Functional Suitability & Functional Completeness & Verifies that the system covers all designated tasks without omission: applicant registration, two-stage document submission and Conditional Revert response, staff verification, and committee review and reallocation. \\
    & Functional Correctness & Assesses, through the Lifecycle Test Case Matrix, whether status transitions follow the Application Lifecycle Path and whether every status change carries a recorded human actor. \\
    & Functional Appropriateness & Examines whether the applicant dashboard and Committee Portal facilitate rapid task completion within the University Scholarship Office's operational context. \\
    Performance Efficiency & Time Behavior & Measures system response times, document upload speed, committee filtering latency, and notification trigger-to-delivery latency. \\
    & Resource Utilization & Monitors cloud resource use, ensuring compound indexing minimizes database load during peak submission windows. \\
    \bottomrule
  \end{tabular}
\end{table}

Evaluation will be conducted through task-based usability testing with a purposive sample of 10 to 15 participants, complemented by post-session instruments: the standardized 10-item System Usability Scale (SUS), which addresses Research Question 4; the Post-Prototype Perceived Effort Survey (Appendix~\ref{app:post-survey}), compared against the pre-study baseline (Appendix~\ref{app:pre-survey}) to address Research Questions 1 and 3; and the Committee Portal Usefulness Survey (Appendix~\ref{app:committee-survey}), which addresses Research Question 2(b). By measuring error rates and time-on-task against the ISO/IEC 25010 dimensions, and by comparing pre- and post-prototype perceived effort for document submission, correction, status tracking, and deadline awareness, ScholarPath AdDU will be empirically validated as a functional, efficient, and user-centric solution to AdDU's scholarship admission bottlenecks.
